\subsection{Integration with respect to a sum of positive measures}

Throughout this No., X denotes a locally compact space, \((\lambda_{\alpha})_{\alpha\in A}\) a summable family of positive measures on X, and \(\nu\) the measure \(\sum_{\alpha\in A}\lambda_{\alpha}\) .

PROPOSITION 1. --- Let \(f\) be a positive numerical function defined on \(X\) . Then

\[
\nu^ {\bullet} (f) = \sum_ {\alpha \in A} \lambda_ {\alpha} ^ {\bullet} (f).\tag{3}
\]

This follows at once from Remark 1, Prop. 11 of §1, No.~3 and Prop. 3 of §1, No.~1.

COROLLARY 1. --- For every compact (resp. open and relatively compact) subset M of X,

\[
\nu (\mathrm{M}) = \sum_ {\alpha \in \mathrm{A}} \lambda_ {\alpha} (\mathrm{M}).
\]

COROLLARY 2. --- For a subset N of X to be locally \(\nu\) -negligible, it is necessary and sufficient that, for every \(\alpha \in A\) , N be locally \(\lambda_{\alpha}\) -negligible.

COROLLARY 3. --- For every function \(f \in \mathcal{F}_{+}(X)\) ,

\[
\nu^ {*} (f) \geqslant \sum_ {\alpha \in A} \lambda_ {\alpha} ^ {*} (f).\tag{4}
\]

The inequality is obvious if \(f\) is not \(\nu\) -moderated, because then \(\nu^{*}(f) = +\infty\) (§1, No.~2, Prop. 7). If \(f\) is \(\nu\) -moderated then \(f\) is \(\lambda_{\alpha}\) -moderated for every \(\alpha \in A\) , because every \(\nu\) -integrable open set is \(\lambda_{\alpha}\) -integrable; the relation (4) then follows at once from (3) and from Prop. 7 of §1, No.~2.

\pandocbounded{\includegraphics[keepaspectratio]{images/part002_632746ad.jpg}}

It can happen that the two members of (4) are not equal, even when A is countable and each of the \(\lambda_{\alpha}\) is a point measure ( \(\S 1\) , Exer. 4 \(a\) )).

PROPOSITION 2. --- Let f be a mapping of X into a topological space G. For f to be \(\nu\) -measurable, it is necessary and sufficient that f be \(\lambda_{\alpha}\) -measurable for every \(\alpha \in A\) .

This follows at once from Cor. 2 of Prop. 11 of §1.

PROPOSITION 3. --- For a mapping f of X into a Banach space F to be essentially \(\nu\) -integrable, it is necessary and sufficient that f be essentially \(\lambda_{\alpha}\) -integrable for every \(\alpha \in A\) and that

\[
\sum_ {\alpha \in \mathrm{A}} \int^ {\bullet} | \mathbf {f} | d \lambda_ {\alpha} <   + \infty .\tag{5}
\]

The family \((\int \mathbf{f}d\lambda_{\alpha})_{\alpha \in \mathrm{A}}\) is then absolutely summable in \(\mathrm{F}\) , and

\[
\int \mathbf {f} d \nu = \sum_ {\alpha \in \mathrm{A}} \int \mathbf {f} d \lambda_ {\alpha}.\tag{6}
\]

Indeed, for \(\mathbf{f}\) to be essentially \(\nu\) -integrable (resp. essentially \(\lambda_{\alpha}\) -integrable), it is necessary and sufficient that \(\mathbf{f}\) be measurable for the measure \(\nu\) (resp. \(\lambda_{\alpha}\) ) and that \(\nu^{\bullet}(|\mathbf{f}|) < +\infty\) (resp. \(\lambda_{\alpha}^{\bullet}(|\mathbf{f}|) < +\infty\) ), by virtue of Prop. 9 of §1, No.~3. The first part of the statement therefore follows at once from Props. 2 and 1. If \(\mathbf{f}\) is essentially \(\nu\) -integrable, the inequality

\[
\sum_ {\alpha \in \mathrm{A}} \left| \int \mathbf {f} d \lambda_ {\alpha} \right| \leqslant \sum_ {\alpha \in \mathrm{A}} \int | \mathbf {f} | d \lambda_ {\alpha} = \nu (| \mathbf {f} |)
\]

implies that the family \(\left(\int\mathbf{f}d\lambda_{\alpha}\right)\) is absolutely summable in F, and that the norm of the sum is less than or equal to the norm of f in \(\overline{\mathcal{L}}_{\mathrm{F}}^{1}(\nu)\) . The set of \(\mathbf{f}\in\mathcal{L}_{\mathrm{F}}^{1}(\nu)\) that satisfy (6) is thus a closed linear subspace H of \(\mathcal{L}_{\mathrm{F}}^{1}(\nu)\) ; now, this subspace is also dense in \(\mathcal{L}_{\mathrm{F}}^{1}(\nu)\) , because it contains the functions of the form \(f\cdot a\) , where \(a\in F\) and f denotes a finite integrable positive function (Prop. 1). Therefore \(\mathcal{H}=\mathcal{L}_{\mathrm{F}}^{1}(\nu)\) and the proposition is established.

Prop. 3 can also be deduced from the general theorem on integration that will be proved in §3 (No.~3, Th. 1).

COROLLARY 1. --- Suppose that f is \(\nu\) -integrable; then f is \(\lambda_{\alpha}\) -integrable for every \(\alpha \in A\) , and formula (6) holds. Conversely, if the set A is finite and f is \(\lambda_{\alpha}\) -integrable for every \(\alpha \in A\) , then the function f is \(\nu\) -integrable.

If \(\mathbf{f}\) is \(\nu\) -integrable, then \(\mathbf{f}\) is essentially \(\nu\) -integrable and \(\nu\) -moderated (§1, No.~3, Cor. of Prop. 9); \(\mathbf{f}\) is therefore essentially \(\lambda_{\alpha}\) -integrable and \(\lambda_{\alpha}\) -moderated, hence \(\lambda_{\alpha}\) -integrable, for every \(\alpha \in A\) . Conversely, if \(A\) is finite and if \(\mathbf{f}\) is \(\lambda_{\alpha}\) -integrable for all \(\alpha \in A\) , then \(\mathbf{f}\) is essentially \(\nu\) -integrable by Prop. 3, and it suffices to verify that \(\nu^{*}(|\mathbf{f}|) < +\infty\) ; this follows at once from the relation \(\nu^{*} = \sum_{\alpha \in A} \lambda_{\alpha}^{*}\) (Ch. IV, §1, No.~3, Prop. 15).

COROLLARY 2. --- Let \(\theta\) be a complex measure on X; set \(\theta_1 = (\mathcal{R}\theta)^+\) , \(\theta_2 = (\mathcal{R}\theta)^-\) , \(\theta_3 = (\mathcal{I}\theta)^+\) , \(\theta_4 = (\mathcal{I}\theta)^-\) . In order that a mapping \(f\) of X into a topological space G (resp. into a Banach space F) be measurable (resp. essentially integrable, integrable) for the measure \(\theta\) , it is necessary and sufficient that it be measurable (resp. essentially integrable, integrable) for each of the measures \(\theta_i\) ( \(i = 1,2,3,4\) ).

If f is measurable (resp. essentially integrable, integrable) for \(\theta\) , then f is by definition measurable (resp. essentially integrable, integrable) for the measure \(|\theta|\) , hence also for the measures \(\theta_{i}\) , which are \(\leqslant |\theta|\) . Conversely, if f is measurable (resp. essentially integrable, integrable) for the measures \(\theta_{i}\) , then Prop. 2 (resp. Prop. 3, Cor. 1 of Prop. 3) implies that f is measurable (resp. essentially integrable, integrable) for the measure \(\theta_{1} + \theta_{2} + \theta_{3} + \theta_{4}\) , which is \(\geqslant |\theta|\) .

